\documentclass[11pt]{article}

\usepackage{amsmath,amssymb}
\usepackage{geometry}
\geometry{margin=1in}

\title{Asymmetric Attractor Structures as a Source of Persistent Matter--Antimatter Imbalance}
\author{}
\date{}

\begin{document}
\maketitle

\begin{abstract}
Recent work has established that discrete, symmetric dynamical systems may
exhibit intrinsically asymmetric attractor structures without explicit
symmetry breaking. In this article, we explore a possible physical
interpretation of such structures in the context of matter--antimatter
imbalance. Without introducing new microscopic dynamics, we argue that
persistent asymmetries may arise from stability properties of the underlying
state space itself. The discussion is deliberately interpretative and does not
modify or replace standard baryogenesis mechanisms.
\end{abstract}

\section{Introduction}

One of the most persistent open problems in fundamental physics is the observed
imbalance between matter and antimatter. Standard quantum field theory permits
the creation of matter and antimatter in symmetric pairs, yet the observable
universe is overwhelmingly matter-dominated.

Traditional explanations rely on dynamical symmetry-breaking mechanisms, often
summarized by the Sakharov conditions. In contrast, the present work explores a
complementary viewpoint: that long-lived asymmetries may originate from the
\emph{stability structure} of the state space itself rather than from explicit
microscopic violations.

This article does not propose a new particle-physics model. Instead, it offers
an interpretation of recently established mathematical results on asymmetric
attractors in symmetric systems.

\section{Mathematical Background}

In a companion paper, a discrete dynamical system was introduced whose evolution
is formally symmetric under phase inversion. Despite this symmetry, the system
exhibits an asymmetric attractor landscape: certain coherent states possess
open basins of attraction, while their symmetry-related counterparts reside on
unstable separatrices.

Crucially, the asymmetry arises at the level of invariant measures rather than
at the level of the evolution rule itself. The dynamics admits symmetric
solutions, but these solutions are dynamically unstable.

\section{Interpretative Mapping}

We now outline a minimal interpretative correspondence:

\begin{itemize}
\item Stable attractors correspond to long-lived, persistent states.
\item Unstable separatrices correspond to transient or short-lived states.
\item Symmetry-related but dynamically inequivalent states naturally lead to
asymmetric occupation probabilities.
\end{itemize}

Within this interpretation, matter and antimatter need not differ in their
fundamental equations of motion. Instead, their effective abundances may differ
because one class of states occupies stable regions of the global phase space,
while the other does not.

\section{Relation to Matter--Antimatter Asymmetry}

From this viewpoint, an initially symmetric ensemble evolves toward a
distribution dominated by dynamically stable states. Even an infinitesimal
bias—arising from fluctuations or initial conditions—becomes amplified by the
geometry of the attractor structure.

Importantly, this mechanism does not violate CPT symmetry at the microscopic
level. The asymmetry emerges statistically and dynamically, not algebraically.

This perspective does not replace baryogenesis scenarios but may complement
them by explaining why extremely small initial asymmetries can persist and
become cosmologically relevant.

\section{Comparison with Standard Approaches}

Conventional approaches attribute matter--antimatter imbalance to explicit
CP-violating interactions combined with out-of-equilibrium dynamics. The
attractor-based interpretation shifts emphasis from interaction terms to global
state-space structure.

Both approaches may coexist. In particular, microscopic CP violation may provide
the initial perturbation, while attractor asymmetry determines its long-term
survival.

\section{Limitations}

The present discussion is intentionally conservative. No claim is made that the
described mechanism quantitatively accounts for the observed baryon asymmetry.
No new experimental signatures are proposed.

The interpretation relies on analogy rather than derivation, and further work
would be required to embed such attractor structures into concrete quantum or
statistical field-theoretic frameworks.

\section{Conclusion}

We have outlined a possible interpretative bridge between asymmetric attractor
structures in symmetric dynamical systems and the persistence of matter over
antimatter. The key insight is that asymmetry may arise from stability properties
of state space rather than from explicit symmetry breaking.

This viewpoint suggests that the geometry of dynamical systems may play a more
fundamental role in cosmological asymmetries than is usually assumed.

\end{document}
